% ============================================================
%  Chapitre 4 — Logique des Défauts
% ============================================================
\chapter{Logique des Défauts}

\section{Fondements théoriques}

\subsection{Motivation}

La logique classique est \textbf{monotone} : ajouter un fait à une base de
connaissances ne peut qu'augmenter l'ensemble des conclusions.
Or, le raisonnement humain est fréquemment \textbf{non-monotone} : on tire des
conclusions par défaut qui peuvent être invalidées par de nouvelles informations.

\begin{tcolorbox}[notebox, title={Exemple de non-monotonie}]
\textit{Par défaut, un artiste est libre de droits.
Mais si l'on apprend qu'il a signé avec un label, cette conclusion est
immédiatement révisée.}
\end{tcolorbox}

\subsection{Théorie des défauts de Reiter (1980)}

Une \textbf{théorie des défauts} est un couple $\Delta = (W, D)$ où :
\begin{itemize}
  \item $W$ est un ensemble de formules du premier ordre (\emph{faits certains}) ;
  \item $D$ est un ensemble de \emph{défauts}.
\end{itemize}

\noindent Un \textbf{défaut} a la forme :
\[
  \frac{\alpha(x) : \beta(x)}{\gamma(x)}
\]
où $\alpha$ est le \emph{prérequis} (doit être dérivable),
$\beta$ est la \emph{justification} (doit être cohérente),
$\gamma$ est la \emph{conclusion}.

\subsection{Extensions}

Une \textbf{extension} $E$ de $\Delta = (W, D)$ est un ensemble maximal et
cohérent de croyances obtenu en appliquant les défauts applicables.
Formellement, $E$ est le plus petit ensemble tel que :
\begin{enumerate}
  \item $W \subseteq E$ ;
  \item $E$ est clos par déduction classique ;
  \item Pour tout défaut $\frac{\alpha:\beta}{\gamma} \in D$,
        si $\alpha \in E$ et $\lnot\beta \notin E$, alors $\gamma \in E$.
\end{enumerate}

\noindent Lorsque deux défauts produisent des conclusions contradictoires,
il existe \textbf{plusieurs extensions} — chacune correspondant à un choix
cohérent de défauts appliqués.

\noindent On distingue :
\begin{itemize}
  \item La \textbf{sémantique crédule} : une conclusion est acceptée si elle
        appartient à \emph{au moins une} extension.
  \item La \textbf{sémantique sceptique} : une conclusion est acceptée si elle
        appartient à \emph{toutes} les extensions.
\end{itemize}

\section{Application — Industrie musicale}

\subsection{Théorie $\Delta = (W, D)$}

\begin{tcolorbox}[defbox, title={Défauts $D$ — domaine musical}]
\begin{align*}
d_1 &: \frac{\text{EstArtiste}(x) : \text{LibreDroits}(x)}{\text{LibreDroits}(x)} \\[4pt]
d_2 &: \frac{\text{SigneLabel}(x) : \lnot\text{LibreDroits}(x)}{\lnot\text{LibreDroits}(x)} \\[4pt]
d_3 &: \frac{\text{EstArtiste}(x) \land \text{AlbumSorti}(x) : \text{Royalties}(x)}{\text{Royalties}(x)} \\[4pt]
d_4 &: \frac{\text{SigneLabel}(x) : \text{PartageRoyalties}(x)}{\text{PartageRoyalties}(x)} \\[4pt]
d_5 &: \frac{\text{IndependantLongtemps}(x) : \text{ExperienceSolo}(x)}{\text{ExperienceSolo}(x)} \\[4pt]
d_6 &: \frac{\text{NouveauContrat}(x) : \lnot\text{LibreDroits}(x)}{\lnot\text{LibreDroits}(x)} \\[4pt]
d_7 &: \frac{\text{PlusDeDeuxAns}(x) \land \text{AlbumSorti}(x) : \text{AlbumCertifie}(x)}{\text{AlbumCertifie}(x)}
\end{align*}
\end{tcolorbox}

\subsection{Résultats par artiste}

\begin{table}[H]
\centering
\caption{Extensions calculées par artiste — domaine musical}
\label{tab:defaut-music}
\begin{tabularx}{\linewidth}{lXll}
\toprule
\textbf{Artiste} & \textbf{Faits $W$} & \textbf{Ext.} & \textbf{Conflit} \\
\midrule
Soolking & EstArtiste + SigneLabel + AlbumSorti & 2 & $d_1$ vs $d_2$ \\
Feu      & EstArtiste + IndependantLongtemps + NouveauContrat & 2 & $d_1$ vs $d_6$ \\
Tif      & EstArtiste + AlbumSorti + PlusDeDeuxAns & 1 & — \\
Flenn (avant) & EstArtiste + AlbumSorti + PlusDeDeuxAns & 1 & — \\
Flenn (après) & + SigneLabel & 2 & $d_1$ vs $d_2$ \\
\bottomrule
\end{tabularx}
\end{table}

\noindent\textbf{Tif — extension unique :}
$E_{\text{tif}} = \{\text{EstArtiste, AlbumSorti, PlusDeDeuxAns,
LibreDroits, Royalties, AlbumCertifie}\}$

\medskip
\noindent\textbf{Soolking — deux extensions (conflit $d_1$ vs $d_2$) :}
\begin{align*}
E_1 &= \{\ldots,\; \text{LibreDroits,\; Royalties,\; PartageRoyalties}\} \\
E_2 &= \{\ldots,\; \lnot\text{LibreDroits,\; Royalties,\; PartageRoyalties}\}
\end{align*}

\noindent Les conclusions \textbf{Royalties} et \textbf{PartageRoyalties}
(issues de $d_3$ et $d_4$) sont dans \textit{toutes} les extensions :
elles sont acceptées en sémantique sceptique.
\textbf{LibreDroits} ne l'est pas : sémantique crédule seulement.

\subsection{Démonstration de non-monotonie}

\begin{tcolorbox}[resultbox, title={Non-monotonie — Flenn}]
\textbf{AVANT} ajout de $\text{SigneLabel(flenn)}$ dans $W$ :
\[
  \text{LibreDroits(flenn)} \in E \quad (\text{unique extension, } d_1 \text{ s'applique})
\]
\textbf{APRÈS} ajout de $\text{SigneLabel(flenn)}$ dans $W$ :
\[
  \text{LibreDroits(flenn)} \notin \bigcap E_i \quad (d_1 \text{ et } d_2 \text{ en conflit,
  2 extensions})
\]
\textbf{Conclusion :} l'enrichissement de $W$ a invalidé une conclusion précédente.
C'est la \textbf{non-monotonie} caractéristique de la logique des défauts.
\end{tcolorbox}

\section{Application — Système judiciaire}

\subsection{Théorie $\Delta = (W, D)$}

\begin{tcolorbox}[defbox, title={Défauts $D$ — domaine judiciaire}]
\begin{align*}
d_1 &: \frac{\text{Accuse}(x) : \text{Innocent}(x)}{\text{Innocent}(x)} \\[4pt]
d_2 &: \frac{\text{Recidiviste}(x) : \text{PeineLourde}(x)}{\text{PeineLourde}(x)} \\[4pt]
d_3 &: \frac{\text{Mineur}(x) : \text{TribunalEnfants}(x)}{\text{TribunalEnfants}(x)} \\[4pt]
d_4 &: \frac{\text{CrimeGrave}(x) : \lnot\text{TribunalEnfants}(x)}{\lnot\text{TribunalEnfants}(x)} \\[4pt]
d_5 &: \frac{\text{Temoignage}(x) : \text{Acquittement}(x)}{\text{Acquittement}(x)} \\[4pt]
d_6 &: \frac{\text{Recidiviste}(x) : \lnot\text{Innocent}(x)}{\lnot\text{Innocent}(x)}
\end{align*}
\end{tcolorbox}

\subsection{Résultats par individu}

\begin{table}[H]
\centering
\caption{Extensions calculées — domaine judiciaire}
\begin{tabularx}{\linewidth}{lXll}
\toprule
\textbf{Individu} & \textbf{Faits $W$} & \textbf{Ext.} & \textbf{Conflit} \\
\midrule
Ali   & Accuse + Temoignage & 1 & — \\
Omar  & Accuse + Recidiviste & 2 & $d_1$ vs $d_6$ \\
Karim & Accuse + Mineur + CrimeGrave & 2 & $d_3$ vs $d_4$ \\
\bottomrule
\end{tabularx}
\end{table}

\noindent\textbf{Ali — extension unique :}
$E_{\text{ali}} = \{\text{Accuse, Temoignage, Innocent, Acquittement}\}$
\quad ($d_1$ et $d_5$ s'appliquent sans conflit)

\medskip
\noindent\textbf{Karim — deux extensions :}
$d_3$ conclut TribunalEnfants, $d_4$ conclut $\lnot$TribunalEnfants.
\textbf{Innocent} (issu de $d_1$) est dans les deux extensions : conclusion sceptique.

\section{Implémentation avec TweetyProject}

\begin{lstlisting}[caption={Calcul des extensions avec Tweety RDL}]
String theory =
    "Art = {soolking}\n"
  + "type(EstArtiste(Art))\n"
  + "type(LibreDroits(Art))\n"
  + "type(SigneLabel(Art))\n"
  + "EstArtiste(soolking)\n"
  + "SigneLabel(soolking)\n"
  // d1 : EstArtiste(X) :: LibreDroits(X) / LibreDroits(X)
  + "EstArtiste(X)::LibreDroits(X)/LibreDroits(X)\n"
  // d2 : SigneLabel(X) :: !LibreDroits(X) / !LibreDroits(X)
  + "SigneLabel(X)::!LibreDroits(X)/!LibreDroits(X)\n";

RdlParser parser = new RdlParser();
DefaultTheory dt = parser.parseBeliefBase(theory);
Collection<Extension> exts =
    new SimpleDefaultReasoner().getModels(dt.ground());
// -> 2 extensions (conflit d1 vs d2)
\end{lstlisting}

\begin{tcolorbox}[notebox, title={Point d'implémentation important}]
Les théories sont \textbf{séparées par individu} (une théorie pour soolking,
une pour feu, etc.) pour éviter l'explosion combinatoire du
\code{DefaultProcessTree} lors du grounding.
La non-monotonie est démontrée en construisant deux théories : l'originale et
l'originale augmentée d'un fait supplémentaire.
\end{tcolorbox}
